\documentclass{article}
\begin{document}
We simulate 400 noise realisations of 250 segments each.
In total we analyse 90000 segments.
Of the 400 realisations, 12.5\% are rejected, which leaves 350.
The bias of the estimator is
\begin{equation}
  b = \frac{c}{N} \label{eq:bias}
\end{equation}
and multiplying by the number of segments gives
\begin{equation}
  b N = c \label{eq:product}
\end{equation}
For the phase we use
\begin{equation}
  \arctan x + \arctan(1/x) = \pi/2 \label{eq:phase}
\end{equation}
Smith reports a bias below 1\% in comparable settings~\cite{smith2020}.
Jones finds that the bias grows with the segment length~\cite{jones2021}.
The estimator is therefore consistent.
\begin{thebibliography}{2}
\bibitem{smith2020} A. Smith, Estimators of noise, J. Stat. 12 (2020) 1.
\bibitem{jones2021} B. Jones, Long segments, arXiv:2101.01234.
\end{thebibliography}
\end{document}
